\documentclass[aspectratio=169]{beamer}

\usepackage[utf8]{inputenc}
\usepackage[brazil]{babel}
\usepackage{graphicx}
\usepackage{hyperref}

% Tema simples e limpo
\usetheme{Madrid}
\usecolortheme{default}
\setbeamertemplate{navigation symbols}{}

% -------------------------------------------------
% INFORMAÇÕES DO EP (OS ALUNOS DEVEM EDITAR)
% -------------------------------------------------

\title{MAC0209 --- Modelagem e Simulação}
\subtitle{Exercício-Programa 1 --- \textit{Floats} e Precisão}
\author[Grupo Plutão]{Grupo Plutão \\[4pt]
Guilherme França Carvalho \\
Ismália Alves Santos Ziemer \\
João Victor Bispo do Nascimento \\
Leonardo Heidi Almeida Murakami \\
Luis Renato Ferraz Natil Martins \\
Murilo Freitas Yonashiro Coelho \\
Vítor Gonçalves Serra \\
Vitor Wallace de Moraes Alves }
\institute[BCC IME-USP]{Bacharelado em Ciência da Computação \\ IME -- USP}
\date{\today}

% -------------------------------------------------

\begin{document}

% =================================================
% SLIDE 1
% =================================================
\begin{frame}
\titlepage
\end{frame}

\begin{frame}{Contexto e Objetivo do EP}

\textbf{Contexto do curso}

\begin{itemize}
    \item Primeira aula, foi apresentada a inicial do curso e alguns dos assuntos que serão tratados.
    \item Erros de precisão em cálculos com float podem levar a resultados inesperados, o que \\deve ser tratado para evitar erros na saída da simulação.
\end{itemize}

\vspace{0.5cm}

\textbf{Objetivo do EP}

\begin{itemize}
    \item Esse EP teve como foco verificar o funcionamento das váriaveis do formato "float" na computação e seus níveis de precisão através de testes feitos na linguagem de programação Python.
\end{itemize}

\end{frame}

% =================================================
% SLIDE 2
% =================================================
\begin{frame}{Trabalho Realizado}

\textbf{Conceitos e implementação}

\begin{itemize}
    \item Representação de números reais em computadores (float seguindo a IEEE-754).
    \item Foram realizados testes para verificar a precisão dos números de ponto flutuante e o resultado de operações realizadas com tais números.
\end{itemize}

\end{frame}

% =================================================
% SLIDE 3
% =================================================
\begin{frame}{Resultados e Conclusões}

\textbf{Resultados}

\begin{itemize}
    \item Principais resultados obtidos.
    \item Interpretação do comportamento observado.
\end{itemize}

\end{frame}

\end{document}
